\subsection*{Appendix B: Accelerators Beyond the GPU}

\paragraph{Exercise B.1}
$x_{mk}$ enters PE$(k,0)$ at cycle $m+k$ and reaches PE$(k,n)$ at $m+k+n$. The partial sum for
$y_{mn}$ starts at PE$(0,n)$ at cycle $m+n$, moves down one row per cycle and is at PE$(k,n)$ at
$m+n+k$, the same cycle as $x_{mk}$, so the MAC happens. It leaves the bottom row at $m+n+K-1$. The
last output ($m=M-1$, $n=N-1$) leaves at cycle $M+N+K-3$, counting from 0, so the pass takes
$M+N+K-2$ cycles. Loading the weights one row per cycle without overlap adds $K$ cycles:
$M+N+2K-2$. \code{simulate\_weight\_stationary(..., include\_weight\_load=True)} reports exactly that.

\paragraph{Exercise B.2}
$K=4096$ and $N=12288$ are multiples of both array sizes, so there is no padding and
$U = M/(M+R+C-2)$. Half utilisation needs $M=R+C-2$: 254 rows on $128\times128$ and 510 on
$256\times256$. Larger arrays need larger batches. TPU serving systems therefore aim for decode
batches of hundreds of sequences per chip (or per TP group), and treat small-batch latency-critical
traffic as a separate, less efficient regime.

\paragraph{Exercise B.3}
Every Qwen3-0.6B dimension (1024, 2048, 3072, and the vocabulary $151{,}936 = 1187\times128$) is a
multiple of 128, so the padding efficiency is exactly 1 for both layouts. GPT-2's 50{,}257 is prime-ish:
the LM head needs $\lceil 50257/128\rceil = 393$ column tiles ($50{,}304$ columns), an efficiency of
0.99907, and $\lceil50257/32\rceil=1571$ tiles of 32 (50{,}272), an efficiency of 0.9997. Padding the vocabulary
to $50{,}304 = 393\times128$ makes both exact. The 47 extra rows are never targets in training, so
their logits are trained towards very negative values or simply sliced off before sampling. They
cost nothing because the hardware was computing padded tiles anyway.

\paragraph{Exercise B.4}
Reduce-scatter along $x$ (16 chips, bidirectional ring, $W$ each way): $\frac{15}{16}\frac{S}{2W}$.
All-reduce along $y$ of the resulting $S/16$: $2\cdot\frac{15}{16}\frac{S/16}{2W}$. All-gather along
$x$: $\frac{15}{16}\frac{S}{2W}$. The total is $\approx0.996\,S/W$, exactly the bandwidth term of one
bidirectional ring over all 256 chips ($2\cdot\frac{255}{256}\frac{S}{2W}$). The difference is latency:
$2\cdot15\cdot2$ ring steps instead of $2\cdot255$. The 2-D schedule can also split $S$ in half and run
$x$-then-$y$ on one half and $y$-then-$x$ on the other concurrently, using both axes' links at once and
approaching $0.5\,S/W$, which a single ring cannot do.

\paragraph{Exercise B.5}
For square blocks, Cannon moves $(p-1)p^2(a+b)$ link-bytes in its shift steps plus the initial skew,
while SUMMA with multicast moves $p^2(p-1)(a+b)$. Cannon therefore costs SUMMA's traffic \emph{plus} the
skew, which is $p\sum_i\min(i,p-i)$ block-hops per operand. For $p=4$ on the test matrices that is 1536
against 1152 bytes. Without multicast SUMMA's line broadcasts cost $p(p-1)/2$ hops each and it becomes
more expensive than Cannon for $p>2$. Cannon is still attractive on a real NoC because every
transfer is a single nearest-neighbour hop on the torus. There are no long multi-hop routes that
contend for links, per-step SRAM use is constant (one $A$ and one $B$ block in flight), and the
traffic is perfectly balanced across links.

\paragraph{Exercise B.6}
A Cannon step on $b\times b$ blocks costs $2b^3/C_c$ seconds of compute against moving two blocks,
$2b^2/W_n$ seconds at 1 byte per element, so $b \ge C_c/W_n = 4\times10^{12}/6.4\times10^{10} = 62.5$:
take $b=64$, a multiple of the 32-element tile. SRAM: two 8-bit input blocks of $64\times64$ are 4\,KB
each, 16\,KB double-buffered, plus a $64\times64$ fp32 accumulator (16\,KB, or 32\,KB
double-buffered), under 50\,KB in total, so it fits easily. The SRAM budget would allow $b$ in the
low hundreds, which leaves room to raise $b$ for slack against NoC contention.

\paragraph{Exercise B.7}
With weights resident in SRAM, a batch-1 token needs $2\times3\times10^8 = 6\times10^8$ FLOPs.
At the card's 8-bit peak (about 466\,TFLOP/s) that is under $\SI{2}{\micro\second}$, so compute is
not the bound. Instead, the per-layer critical path is: every layer's ops run in sequence, each with
dispatch and synchronisation overheads of microseconds, plus NoC transfers of activations between
the cores that hold consecutive layers. For a model of a few dozen layers with several ops per
layer, those overheads put the realistic ceiling in the low thousands of tokens per second. The
DRAM-bound alternative is $576\times10^9/3\times10^8\approx1900$ tokens/s. SRAM residency helps
most when it removes DRAM from the loop \emph{and} the op overheads are amortised: through batching,
op fusion, or a static schedule that pipelines layers across cores.

\paragraph{Exercise B.8}
(a) 1000 concurrent users of an 8B model is a throughput and KV-capacity problem: a high-bandwidth,
high-capacity machine (H100- or TPU v5p-class) with continuous batching and paged KV caches.
(b) One developer running a 70B model locally needs capacity: a 128\,GB unified-memory machine (GB10)
holds a 4-bit 70B model with a long cache, at single-digit tokens/s. (c) GRPO on a 1.7B model needs
fast rollouts and training in one software ecosystem, which today means a CUDA GPU node, or TPUs with a
JAX RL stack. (d) A latency-critical single-user 8B assistant is best served by an SRAM-resident
design (Groq, Cerebras), or a high-bandwidth GPU with speculative decoding.

\paragraph{Exercise B.9}
A sequence of 4096 tokens has $0.604$\,GB of bf16 KV cache. KV traffic equals weight traffic when
$B\cdot0.604 = 16.4$ (bf16 weights), $B\approx27$, or $B\cdot0.604=8.2$ (fp8 weights), $B\approx14$.
Quantising the cache to 8 bits halves the per-sequence term and doubles both crossovers, to about 54 and 27.
With fp8 weights, an 8-bit cache is the natural pairing.

\paragraph{Exercise B.10}
Per-token time is $28.5\,\text{ms}/n$ (8.2\,GB of fp8 weights at 288\,GB/s, split over $n$ chips) plus
$72\cdot2(n-1)\alpha$ of all-reduce latency. With $\alpha=\SI{2}{\micro\second}$: 14.5, 8.0 and
5.6\,ms for $n=2,4,8$ (69, 125 and 179 tokens/s), so TP keeps paying off. With
$\alpha=\SI{20}{\micro\second}$: 17.1, 15.8 and 23.7\,ms (58, 63 and 42 tokens/s), so scaling stops
at $n=4$, and $n=8$ is slower than $n=2$. For many users, run several TP-2 or TP-4 groups as data-parallel
replicas and batch requests within each. Batching amortises both the weight streaming and the
per-step collective latency over many tokens.
